\documentclass{beamer}
%\usepackage[T1]{fontenc}
\usepackage[utf8]{inputenc}
\usepackage{booktabs}
\usepackage{url}
\usepackage{listings}

\newcommand{\Heading}[1]{\hskip1cm\textbf{#1}\\}

\mode<presentation> {
	\usetheme{FZU}
}

%\usepackage[orientation=landscape,size=custom,width=200,height=120,scale=1.9]{beamerposter}
\usepackage[orientation=portrait,size=a0,scale=1.4]{beamerposter}
%\usepackage{times}

\title{\Huge Deska: Tool for Central Administration of a Grid Site}
\author{\textbf{Jan Kundrát}, Martina~Krejčová, Lukáš~Kerpl, Tomáš~Hubík}
\institute[Institute of Physics]{Institute of Physics, AS CR, Prague, Czech
Republic}
\def\Email{kundratj@fzu.cz}

\lstset{
basicstyle=\ttfamily\scriptsize,
xleftmargin=1cm
}

\begin{document}
\begin{frame}[fragile]

\begin{columns}[t]

\begin{column}{.48\linewidth}
\begin{block}{Introduction}
Recent developments in hardware density brought unmatched \alert{challenges} to
the shoulders of system administrators.  Number of CPUs and machines to be
managed \alert{grows rapidly}, yet the number of staffers tends to remain
\alert{stagnant}.  Some attempts at making the infrastructure work reliably,
like adding more monitoring tools, could actually make people's life harder due
to the need to \alert{duplicate information} about hosts even more.
\end{block}

\begin{block}{The Duplicated Information}
Let's check where data for each machine deployed in the network have to be stored
and maintained:
\begin{columns}[T]
\begin{column}{.4\linewidth}
\begin{itemize}
    \item HW inventory DB
    \item Warranty \& issue tracking
    \item Switch port configuration
    \item DHCP server
    \item DNS
    \item Cfengine roles
\end{itemize}
\end{column}
\begin{column}{.5\linewidth}
\begin{itemize}
    \item Torque's CPU multipliers
    \item MRTG \& RRD network graphs
    \item Nagios
    \item Ganglia
    \item Munin
    \item \ldots
\end{itemize}
\end{column}
\end{columns}
\vspace{1em}

Clearly, we should manage that information from one \alert{central place}.
\end{block}

\begin{block}{The CLI approach}

\Heading{Web GUIs}
The institute's staffers have used an in-house web-based tool for HW management.
It has serious limitations like the lack of \alert{bulk operations} which make
working with the application an unpleasant experience.  In contrast, such
operations are best dealt with using \alert{CLI interfaces} and \alert{version
control systems}.

\Heading{The CLI solution}
The idea is to develop a system controlled in a similar way to how one
configures Cisco IOS, ie. using a nice \alert{shell} with friendly
\alert{tab-completion}.

\Heading{Backup and Change Management}
The \alert{textual structure} of the database also allows for extremely easy
backups.  In addition, keeping track of changes, even in an external system, is
trivial.
\end{block}

\begin{block}{Database Structure}
\Heading{List of Objects}
Similar to how the IOS works, we define \alert{lists} of \alert{top-level
objects}.  A good candidate for top-level objects are \alert{machines},
\alert{different HW models}, \alert{racks} etc.

\Heading{Assigning Properties}
For each object, a well-defined set of \alert{properties} or \alert{attributes}
is supported.

For example, for any given \alert{host}, it makes sense to store its \alert{rack
location}, \alert{warranty number}, \alert{HW-specific information} etc.
\end{block}

\begin{block}{Abstract and Concrete Objects}
Intuitively, two kinds of objects should be managed, the \alert{physical}
entities like a {\tt host} or a {\tt rack}, and \alert{abstract} elements like
\alert{hardware model} ({\tt hwmodel}) or a {\tt network}.
\end{block}

\begin{block}{Inheritance and Templates}
It makes a lot of sense to try to re-use existing definitions as much as
possible.  Consider the racks -- most of them are very similar and one can
exchange machines between them freely, yet the racks differ in number of
available bays.  This is how we do that in Deska:

\begin{lstlisting}
boxmodel generic-rack
    outer width 700
    outer depth 1300
    outer height 2000

    bay height 44 order bottom-up start-at 1
end

boxmodel APC-rack
    template generic-rack

    outer height 2300

    bay height 48
end
\end{lstlisting}

In this case, the {\tt APC-rack} is \alert{derived} from the {\tt generic-rack}.
Nothing prevents us from using the {\tt generic-rack} directly, and all
equipment that fits into a {\tt generic-rack} could be put into the {\tt
APC-rack} without any problem -- it's just a \alert{special case} of a rack
after all.
\end{block}

\end{column}

\begin{column}{.48\linewidth}

\begin{block}{Using the Data}
A set of \alert{generators} is used for converting data from the database into
anything from \alert{configuration files} to \alert{switch configuration}.
Using Deska's API, generators retrieve identifiers of objects matching certain
\alert{criteria} and process them appropriately.

This is used for generating DNS \alert{ zone files}, \alert{DHCP} configuration
and even \alert{cfengine} classes.
\end{block}

\begin{block}{Versioning and Atomicity}
\Heading{Sandbox}
Performed changes are not instantly applied to live data.  Each CLI session
lives in a sandboxed environment, and only an explicit \alert{commit} operation
propagates the changes into the production version.

\Heading{Version Control}
Administrators are presented a textual \alert{diff-like} changeset upon each
commit attempt.  The output indicates both \alert{DB changes} and modifications
to the \alert{generated files}.  The intention is to operate as safely as
possible by clearly showing the \alert{impact of the change}.
\end{block}

\begin{block}{Implementation Details}
\begin{columns}[T]
\begin{column}{.4\linewidth}
\Heading{Choice of Tools}
\begin{itemize}
    \item C++, STL, TR1,\ldots
    \item PostgreSQL
    \item SSH and PAM for authentication, authorization and the transport
        channel
\end{itemize}
\end{column}

\begin{column}{.6\linewidth}
\Heading{Target Platform}
\begin{itemize}
    \item The goal is to remain cross-platform; a reasonable target for
        production are both 32bit and 64bit RHEL5 clones
    \item The client should work on any recent Linux desktop distribution.
\end{itemize}
\end{column}
\end{columns}

\end{block}

\begin{block}{Extensibility}
The only \alert{hard-coded assumption} is that the database consists of a list
of objects with well-defined \alert{structure}.  Each kind of objects has a set
of allowed \alert{attributes} of a specified \alert{data type}.

Predefined types of objects are shipped with Deska.  These objects allow a
full and \alert{complete description} of all the equipment used at Prague site.
However, we recognize the need for \alert{runtime configuration} of what is to
be stored in the database, such as keeping detailed track of the Infiniband
infrastructure.

Therefore, it is possible to define new \alert{kinds} of objects by a simple
plugin mechanism, as well as to \alert{extend} the properties of existing
objects.  The database structure is augmented upon plugin activation.
\end{block}

\begin{block}{Defining Plugins}
Trying to be consistent with the ``don't duplicate information'' motto, we are
investigating ways to further \alert{automate} plugin definitions.  The most
likely option to follow is using Postgres' code for extracting table structure
from the \alert{SQL} definitions.  Re-using SQL would allow us to retain the
ability to use native triggers for ensuring \alert{data consistency}, as well as
avoid having to invent yet another constraint syntax.
\end{block}

\begin{block}{Links}
\begin{itemize}
    \item \url{http://projects.flaska.net/projects/show/deska}
    \item \url{http://www.farm.particle.cz/}
    \item Mailing list \url{deska@lists.flaska.net}
\end{itemize}
\includegraphics[width=\linewidth,trim=0 5.8mm 0 0]{deska_logo.pdf}
\end{block}


\end{column}

\end{columns}
\vfill
\end{frame}

\end{document}
